\bookchapter{Bab 10: Pembayaran Bergerak \& Perkhidmatan Hidup}

\chapterintro{Masukkan nombor kad. Sekarang kebanyakan App menyokong terus gambar kenal pasti nombor kad, tak perlu taip satu-satu digit.}

\booksection{Bank akan hantar SMS kod pengesahan ke nombor telefon yang anda daftarkan. Isi, tetapkan satu kata laluan pembayaran 6 digit.}

Kata laluan pembayaran ini, kata laluan kad bank dan kata laluan skrin kunci anda, terbaiknya tiga kata laluan berbeza, jangan mudahkan tetapkan sama.

\begin{enumerate}
  \item Ada tip kecil: jika risau keselamatan, boleh buka akaun elektronik Kategori II atau III dalam App bank, khas untuk mengikat pembayaran telefon. Kad ini hanya letak wang belanja harian, walaupun berlaku masalah, kerugian juga terhad.
  \item 2. Kod pembayaran adalah "siapa lihat siapa kumpul wang"
  \item Banyak orang tak berjaya bezakan kod pembayaran dan kod penerimaan, perbezaan ini sangat penting.
  \item Kod pembayaran
\end{enumerate}

: Kod bar + kod QR yang anda tunjukkan kepada juruwang imbas. Kod ini setara dengan “dompet elektronik” anda, sesiapa imbas boleh tolong wang dari akaun anda. Jadi

\booksection{jangan sekali-kali tangkap skrin dan hantar kepada sesiapa}

. Kod ini automatik disegarkan semula setiap minit, untuk elakkan digunakan tanpa kebenaran.

\begin{itemize}
  \item \textbf{Kod penerimaan}: Orang imbas kod ini untuk pindah wang kepada anda. Kod ini boleh ditangkap skrin disimpan, boleh dihantar kepada orang lain, selamat.\textbf{Kenal pasti perbezaan ini, boleh elakkan banyak perangkap penipuan. Sesiapa yang suruh anda "tangkap skrin kod pembayaran dan hantar", 100\% penipu.}3. Tetapan keselamatan ini, 5 minit tetapkan boleh jimat banyak masalah
  \item \textbf{Pembayaran cap jari/wajah}: Dalam tetapan pembayaran Alipay dan WeChat, aktifkan pembayaran biometrik. Jadi setiap pembayaran perlu tekan cap jari atau imbas wajah, walaupun telefon hilang, orang yang jumpa juga tak boleh guna wang anda.
\end{itemize}

Kunci keselamatan

\booksection{: Alipay dan WeChat ada fungsi “Kunci Keselamatan”. Selepas diaktifkan, masuk ke halaman “Dompet” App perlukan pengesahan identiti tambahan. Laluan: Alipay Saya Tetapan Akaun \& Keselamatan Pusat Keselamatan Kunci Keselamatan.}

\begin{itemize}
  \item \textbf{Semak pembayaran tanpa kata laluan}: Banyak orang langgan pelbagai pembayaran tanpa kata laluan, sendiri pun lupa. Pergi ke "Tetapan Pembayaran" Alipay dan WeChat, cari "Pembayaran Tanpa Kata Laluan/Deduct Automatik", tengok perkhidmatan apa yang dilanggan, tutup yang tak diperlukan.
  \item \textbf{Pengurusan peranti log masuk}: Selalu semak peranti mana yang pernah log masuk akaun anda. Jika jumpa peranti yang tak dikenali, terus keluarkan. Laluan: Alipay Tetapan Akaun \& Keselamatan Pengurusan Peranti Log Masuk.
  \item \textbf{4. Bas \& Metro: Telefon adalah kad transit}Sekarang naik bas dan metro, sudah ramai tak guna kad transit fizikal. Telefon imbas pintu gerbang, ada dua cara:
  \item \textbf{Kad transit maya dalam dompet telefon}: Dalam Apple Wallet, Huawei Wallet, Xiaomi Wallet terus buka satu kad transit maya. Kelebihan cara ini: semasa imbas tak perlu buka kunci telefon, terus dekatkan sahaja, kelajuan tindak balas paling pantas, dan boleh imbas walaupun tiada rangkaian. Kad transit maya di sesetengah bandar juga ada diskaun pertukaran sama seperti kad fizikal.
\end{itemize}

\booksection{Kod menaiki Alipay/WeChat}

: Cari “Kod Menaiki” dalam App, buka kod menaiki bas atau metro bandar anda. Kelebihan cara ini: tak perlu tambah nilai, terus tolong dari baki atau kad bank yang diikat. Tapi perlukan sambungan internet, dan sebelum imbas perlu buka App cari kod.

\begin{itemize}
  \item \textbf{Saya peribadi lebih cadangkan guna kad transit maya dompet telefon. Walaupun kali pertama perlu tambah 10-20 yuan, tapi guna memang jauh lebih mudah, terutama bila kejar metro, tak perlu berhenti buka kunci telefon cari App.}5. Bayaran bil utiliti, penghantaran makanan, teksi
  \item \textbf{Bayaran air elektrik gas}: Dalam "Bayaran Bil Hidup" Alipay atau WeChat, ikat nombor akaun air, elektrik, gas rumah anda. Selepas itu setiap bulan terima bil, terus klik boleh bayar. Nombor akaun biasanya boleh dijumpai pada bil kertas rumah anda, atau telefon soal khidmat pelanggan.
\end{itemize}

Perbandingan harga makanan penghantaran

\booksection{: Kedai yang sama di Meituan dan Ele.me selalunya harga dan promosi berbeza. Sebelum pes tengok kedua-dua belah, kadang-kadang boleh jimat beberapa yuan. Tambahan pula, fungsi “Pin Hao Fan” Meituan, pesan bersama kawan, harga boleh jauh lebih murah.}

\begin{itemize}
  \item \textbf{Keselamatan teksi}: Dalam DiDi atau Gaode Taxi, ingat tetapkan "Kenalan Kecemasan" dan "Kongsi Perjalanan". Semasa panggil teksi, App akan automatik kongsi perjalanan anda kepada kenalan yang ditetapkan. Tambahan pula, lewat malam panggil teksi utamakan "Teksi" bukan "Kuai Che", teksi ada pengawasan industri, tahap keselamatan lebih tinggi.
  \item \textbf{6. Urus insurans perubatan, keselamatan sosial, kerajaan dengan telefon}Perubahan beberapa tahun ini sangat besar, banyak perkara yang dahulu perlu beratur di dewan, sekarang dalam telefon beberapa minit selesai.
  \item \textbf{Voucher elektronik insurans perubatan}: Cari "Voucher Elektronik Insurans Perubatan" dalam Alipay atau WeChat, selepas diaktifkan, pergi hospital beli ubat terus imbas telefon, tak perlu bawa kad insurans perubatan fizikal. Kali pertama perlukan pengesahan pengenalan wajah, selepas itu setiap kali guna akan jana satu kod dinamik.
\end{itemize}

\booksection{Query keselamatan sosial}

: Muat turun App “Zhangshang 12333” atau guna “Pusat Warganegara” dalam Alipay, boleh semak rekod pembayaran keselamatan sosial, baki akaun pencen, tuntutan elaun pengangguran dll.

\begin{itemize}
  \item \textbf{Cukai pendapatan}: Penjelasan cukai pendapatan tahunan Mac hingga Jun, boleh selesaikan dalam App "Cukai Pendapatan Perseorangan". Ingat tengok ada potongan tambahan khas yang boleh diisytiharkan (pendidikan anak, faedah pinjaman rumah, menyara warga tua dll.), semua ini boleh bantu anda kurang bayar cukai atau dapat bayaran balik cukai.
  \item \textbf{Ingat ini, elak penipuan jaga keselamatan}Pihak berkuasa, bank, khidmat pelanggan, takkan suruh anda pindah wang ke "akaun selamat".
  \item \textbf{Sesiapa suruh anda tangkap skrin kod pembayaran dan hantar, semua penipu.}Jangan klik pautan asing, terutama SMS seperti "ETC tamat tempoh", "Kad keselamatan sosial dihentikan".
\end{itemize}

\begin{tcbwarn}{Bayaran balik belian dalam talian akan kembali ke kaedah asal, tak perlu anda berikan nombor kad bank.}
\begin{itemize}
  \item Buat sekarang: Pemeriksaan keselamatan Apl pembayaran anda
  \item Buka Alipay dan WeChat, sahkan pembayaran cap jari/wajah telah diaktifkan.
  \item Cari fungsi "Kunci Keselamatan" dan aktifkan.
  \item Semak senarai "Pembayaran Tanpa Kata Laluan/Deduct Automatik", tutup perkhidmatan yang tak diperlukan lagi.
\end{itemize}

\end{tcbwarn}

\begin{tcbpractice}{Tengok senarai "Peranti Log Masuk", buang peranti yang tak dikenali.}
\begin{enumerate}
  \item Aktifkan voucher elektronik insurans perubatan anda, kali seterusnya pergi hospital tak perlu bawa kad fizikal.
  \item Bab 11: Fotografi \& Pengurusan Gambar
  \item Sejujurnya, kebanyakan orang ambil gambar dengan telefon, hasilnya jauh kurang berbanding potensi telefon itu sendiri. Bukan teknik anda tak bagus, tapi banyak tetapan kecil dan tabiat kecil yang anda tak perasan. Bab ini tak cakap benda kosong, hanya benda yang anda boleh guna terus.
  \item 1. Adakah anda juga buat kesilapan ini?
  \item Mari cakap beberapa masalah paling biasa, tengok anda kena berapa:
\end{enumerate}

\end{tcbpractice}
